\section{Related Work}
\label{sec:related}

\para{Instruction-following benchmarks.} The evaluation of LLM instruction following has progressed from broad capability assessments to fine-grained constraint checking. IFEval~\cite{zhou2023ifeval} introduced 25 verifiable instruction types (keyword inclusion, format requirements, length constraints) with deterministic verification, establishing the paradigm of judge-free evaluation. However, all 25 types are unconditional---they test whether a model can satisfy a flat constraint, not whether it can conditionally apply one. FollowBench and AdvancedIF~\cite{he2025advancedif} extend this paradigm to multi-level and multi-turn settings but similarly focus on unconditional adherence. Young et al.~\cite{young2025models} evaluate 256 LLMs on 20 instruction-following prompts, revealing consistent failure patterns across model families, though without isolating conditional logic. Our work complements these benchmarks by targeting a specific capability---conditional rule application---that they do not measure.

\para{Compositional and complex instructions.} ComplexBench~\cite{wen2024complexbench} introduces a taxonomy of constraint compositions---And, Chain, and Selection---where Selection directly corresponds to if-then-else branching. Their finding that Selection accuracy (0.765) is lower than And accuracy (0.881) provides early evidence that conditional branching is harder than parallel constraints. However, ComplexBench evaluates Selection as one of many composition types rather than systematically varying the type of condition. TOD-ProcBench~\cite{ghazarian2025todprocbench} operationalizes if-then logic as multi-level condition-action pairs in customer service dialogues, finding that even frontier models achieve only 39.3\% on combined instruction retrieval and action prediction. SIFo~\cite{chen2024sifo} tests causally dependent sequential instructions, where implicit conditionality arises from chained dependencies. We build on these findings by isolating condition type as the independent variable and measuring both false positive and false negative rates---a distinction not made in prior work.

\para{Rule following and logical reasoning.} RuleBERT~\cite{saeed2021rulebert} demonstrates that pre-trained language models can learn soft Horn rules through fine-tuning, achieving strong generalization to novel if-then patterns. However, this work addresses learned rule following through parameter updates rather than in-context rule specification. Liu et al.~\cite{liu2024incomplete} reveal a dissociation between instruction following, few-shot learning, and instruction inference, suggesting that the format of if-then rules (examples versus explicit descriptions) may affect adherence. Zhao et al.~\cite{zhao2024icl} show that in-context learning plateaus after 10--30 examples and underperforms instruction fine-tuning on multi-turn tasks, raising questions about the reliability of in-context rule specification for conditional behavior.

\para{Verification of instruction following.} Recent work has developed sophisticated verification methods for instruction adherence. IF-CRITIC~\cite{wen2025ifcritic} decomposes instructions into constraint checklists for fine-grained evaluation. NSVIF~\cite{su2026nsvif} formulates verification as constraint satisfaction over both logical and semantic constraints. Our approach takes a simpler path: by requiring models to include or exclude distinctive marker strings, we achieve fully deterministic verification without any judge model, at the cost of testing a narrower (but precisely controlled) form of instruction following.
